\documentclass{article}
\usepackage{mathnotes}

\title{Discrete Entropy}
\slug{discrete-entropy}

\begin{document}

\section{Capacity of a Discrete Channel}

\begin{definition}[Discrete Channel]
A system whereby a \@{sequence} of choices from a \@{finite} \@{set} of elementary symbols $S_1, \dots, S_n$ can be transmitted from one point to another. Each of the symbols $S_i$ is assumed to have a certain duration in time $t_i$ seconds (not necessarily the same for different $S_i$).
\end{definition}

\begin{definition}[Capacity]\synonyms{"channel-capacity"}
The \textbf{capacity} $C$ of a \@{discrete-channel} is given by

\[ C = \lim_{t \to \infty} \frac{\log{N(t)}}{t}, \]

where $N(t)$ is the number of allowed \@{signals} of duration $t.$
\end{definition}

\begin{note}
Given an information source where all symbols are of the same time duration, and each symbol represents $s$ bits of information (because it is chosen freely among $2^s$ symbols), and the channel can transmit $n$ symbols per second then the \textbf{\@{capacity}} $C$ of the channel is defined to be $ns$ bits per second.

In the more general case we have to deal with symbols of various lengths that take different amounts of time to transmit, and so capacity measures not the number of symbols transmitted per second but the amount of information transmitted per second, retaining bits per second as its unit.

If $N(t)$ represents the number of \@{sequences} of duration $t,$ then

\[ N(t) = N(t - t_1) + N(t - t_2) + \cdots + N(t - t_n). \]

This number is equal to the sum of the numbers of sequences ending in $S_1, S_2, \dots, S_n.$ This is a recursive definition - if the last symbol is $S_i,$ and then we have $t - t_i$ time remaining for the previous symbols in the sequence, and we repeat the same question on that remainder.

A "well-known result" in finite differences tells us that

\[ N(t) \sim A W_0^t \quad \text{as } t \to \infty, \]

where $A$ is constant and $W_0$ is the largest real solution of the \@{characteristic-equation}:

\[ W^{-t_1} + W^{-t_2} + \cdots + W^{-t_n} = 1, \]

and therefore 

\[ C = \lim_{t \to \infty} \frac{\log{A W_0^t}}{t} = \log{W_0}. \]
\end{note}

\begin{example}
Consider Morse code as used in telegraph. We have the following rules:

\begin{itemize}
\item A \textbf{dot} symbol consists of one time unit of line closure followed by one time unit of line open.
\end{itemize}

\begin{itemize}
\item A \textbf{dash} symbol consists of three time units of line closure followed by one time unit of line open.
\end{itemize}

\begin{itemize}
\item A \textbf{letter space} symbol consists of three time units of line open.
\end{itemize}

\begin{itemize}
\item A \textbf{word space} symbol consists of six time units of line open.
\end{itemize}

We're never allowed to send two space symbols in a row, because two sequential letter space symbols are indistinguishable from a word space symbol.

Now, our possible terminating states are

\begin{tabular}{lll}
Symbol & Composition & Duration (time units) \\
\hline
Dot & 1 unit closed + 1 unit open & 2 \\
Dash & 3 units closed + 1 unit open & 4 \\
Letter space preceded by dot & 1 unit closed + 4 units open & 5 \\
Letter space preceded by dash & 3 units closed + 4 units open & 7 \\
Word space preceded by dot & 1 unit closed + 7 units open & 8 \\
Word space preceded by dash & 3 units closed + 7 units open & 10 \\
\end{tabular}

Note that we had to consider the last two symbols in the case of the last symbol being a space in order to deal with our "no sequential spaces" constraint.

Now we have

\[ N(t) = N(t - 2) + N(t - 4) + N(t - 5) + N(t - 7) + N(t - 8) + N(t - 10). \]

So, our \@{characteristic-equation} is

\[ W^{-2} + W^{-4} + W^{-5} + W^{-7} + W^{-8} + W^{-10} = 1. \tag{a} \]

With a substitution of $w = 1/W$ we get

\[ w^{2} + w^{4} + w^{5} + w^{7} + w^{8} + w^{10} = 1, \]

and $w_0,$ the largest positive root of this equation, found numerically, is about $0.6882,$ and so

\[ C = -\log{(w_0)} \approx 0.539 \text{ bits per unit of time}. \]

Note that in this Morse telegraphy system, we have two states that the channel can be in, based on what the previous symbol transmitted was.

\begin{itemize}
\item If the previous symbol transmitted was a space, we're in state $1,$ and the next symbol can only be a dot or a dash, and the state will change.
\end{itemize}

\begin{itemize}
\item If the previous symbol transmitted was not a space, we're in state $2,$ and the next symbol can be anything, and the state may or may not change depending on what the next symbol is.
\end{itemize}

We can think of these states being the \@{nodes} of a \@{directed-graph} with the transitions between them being the \@{edges}. A generalization of this is given in the theorem below.
\end{example}

\begin{theorem}
Let $b_{ij}^{(s)}$ be the duration of the $s$-th symbol which is allowable in state $i$ and leads to state $j.$ Then the \@{channel-capacity} $C$ is equal to $\log{W}$ where $W$ is the largest real root of the \@{determinantal-equation}

\[ \left | \sum_{s} W^{-b_{ij}^{(s)}} - \delta_{ij}  \right | = 0, \]

where $\delta_{ij}$ is the \@{Kronecker-delta}.
\end{theorem}

\detach

\begin{example}
For example, for our Morse telegraphy example, we have

\[ \begin{vmatrix} -1 & W^{-2} + W^{-4} \\ W^{-3} + W^{-6} & W^{-2} + W^{-4} - 1 \end{vmatrix} = 0. \]

Expanding the \@{determinant} on the LHS gives the \@{characteristic-equation} $(a)$ above.
\end{example}

\section{Information and Entropy in a Discrete Channel}

\begin{definition}[Self-Information]\synonyms{"surprisal", "Shannon information", "Information content"}
Given a \@{real-number} $b > 1$ and an outcome $x$ of a discrete \@{random-variable} $X$ with probability mass function $p,$ the \textbf{self-information} of $x$ is defined as the negative \@{log-probability}

\[ I(x) := - \log_{b}{p(x)}. \]
\end{definition}

\begin{remark}
We can view \@{self-information} as an alternative casting of \@{probability}, like how \@{odds} are, with some desirable properties:

\begin{itemize}
\item An event with probability $1$ provides no information - we knew it was going to happen so it happening tells us nothing new.
\end{itemize}

\begin{itemize}
\item The less probable an event is, the more \@{surprisal} it yields.
\end{itemize}

\begin{itemize}
\item \@{self-information} is \@{additive} - if two individual events are measured separately, the total amount of information is the sum of the \@{self-informations} of the individual events.
\end{itemize}

The function given above is the unique function of probability (up to a multiplicative scaling factor) that satisfies these three properties.
\end{remark}

\begin{remark}
In a related vein we can develop a function that asks how much choice is involved in the selection of an event, or equivalently, how much uncertainty there is in the outcome. If $p_1, p_2, \dots, p_n$ are the probabilities of events occurring, we want a measure $H(p_1, p_2, \dots, p_n)$ that has the following properties:

\begin{enumerate}
\item $H$ should be \@{continuous} in the $p_i.$ That is, a small change in a $p_i$ should result in a small change in $H.$
\end{enumerate}

\begin{enumerate}
\item If all the $p_i$ are equal, $p_i = \frac{1}{n},$ then $H$ should be a \@{monotonic} \@{increasing} \@{function} of $n.$ With equally likely events there is more choice, or uncertainty, when there are more possible events.
\end{enumerate}

\begin{enumerate}
\item If a choice is broken down into two successive choices, the original $H$ should be the weighted sum of the individual values of $H.$ For example, we should have that
\end{enumerate}

\[ H(\frac{1}{2}, \frac{1}{3}, \frac{1}{6}) = H(\frac{1}{2}, \frac{1}{2}) + \frac{1}{2}H(\frac{2}{3}, \frac{1}{3}). \]
\end{remark}

\begin{definition}[Entropy]\synonyms{"uncertainty"}
Given a discrete \@{random-variable} $X,$ which may be any \@{element} $x$ within the set $\mathcal{X},$ and is distributed according to $p: \mathcal{X} \to [0,1],$ the \textbf{entropy} is

\[ H(X) := - \sum_{x \in \mathcal{X}} p(x) \log{p(x)}. \]

We can choose different bases for the logarithm; throughout these notes a bare $\log$ means $\log_2,$ giving the unit of \@{bits}.

An alternative, equivalent definition is that \textbf{entropy} is the \@{expected-value} of the \@{self-information} of a \@{random-variable}:

\[ H(X) = \mathbb{E}\left[I(X) \right ] =  \mathbb{E} \left [ -\log p(X) \right ]. \]

The unit for entropy is bits per symbol.
\end{definition}

\begin{definition}[Commensurable]\synonyms{"commeasurable"}
\@[Real numbers]{real-numbers} $a_1, \dots, a_n$ are \textbf{commensurable} if there exists a common measure $m \in \reals_{>0}$ such that each $a_i$ is an \@{integer} multiple of it:
\[ a_i = k_i\, m, \qquad k_i \in \integers, \quad i = 1, \dots, n. \]
Equivalently, all pairwise ratios are rational: $a_i / a_j \in \rationals$ for all $i, j$ (with $a_j \neq 0$).
\end{definition}

\begin{theorem}
The only $H$ satisfying the three required properties above is the \@{entropy} function defined above, up to multiplication by a constant.
\end{theorem}

\begin{proof}
Assume we have a function $A(n)$ that satisfies the three properties listed above, such that

\[ A(n) = H(\frac{1}{n}, \frac{1}{n}, \dots, \frac{1}{n}). \]

Then, by property (3) above, we can decompose a choice from $s^m$ equally likely possibilities into a sequence of $m$ choices each from $s$ equally likely possibilities. For example, if we have $2^4$ equally likely possibilities, the probability of any given event is $1/16$. If we instead we have a series of $4$ choices each with $1/2$ probability, we end up with $1/16$ as the probability of any specific sequence of events. So, we have that $A(s^m) = m A(s).$

Now, with arbitrarily large $n,$ we can also have $t^n$ such that $A(t^n) = nA(t),$ by the same logic, and we can pick $m$ such that

\[ s^m \leq t^n < s^{m+1}. \]

Now, we can take the logarithm of each term to get

\[ m \log{s} \leq n \log{t} < (m+1)\log{s}, \]

and dividing by $n \log{s}$ gives

\[ \frac{m}{n} \leq \frac{\log{t}}{\log{s}} < \frac{m}{n} + \frac{1}{n}, \]

and because $n$ is arbitrarily large, 

\[ \left | \frac{m}{n} - \frac{\log{t}}{\log{s}} \right | < \epsilon, \]

where $\epsilon$ is arbitrarily small.

By property (2) of $A(n)$ (it is a \@{monotonically-increasing} function of $n$,)  

\bal
A(s^m) & \leq A(t^n) \leq A(s^{m+1}) \\
mA(s) & \leq nA(t) \leq (m+1)A(s).
\eal

Then, dividing by $nA(s)$ gives

\[ \frac{m}{n} \leq \frac{A(t)}{A(s)} \leq \frac{m}{n} + \frac{1}{n} \text{ or } \left | \frac{m}{n} - \frac{A(t)}{A(s)} \right | < \epsilon, \]

Now, by the \@{triangle-inequality}, we have that

\bal
\left | \frac{A(t)}{A(s)} - \frac{\log{t}}{\log{s}} \right | & \leq 2 \epsilon \\
\left | A(t) - \frac{A(s)}{\log{s}} \log{t} \right | & \leq 2 \epsilon A(s).
\eal

Since $\epsilon$ can be arbitrarily small, we have that $A(t) = K \log(t),$ with $K > 0$ so that property (2) holds. Now we know what $A(n)$ is, and thus what $H$ is when we have equal probabilities for all events.

Now let's say that we have a choice from $n$ possible events with \@{commensurable} probabilities $p_i = \frac{n_i}{\sum n_i}.$ We can break down a choice from $\sum n_i$ possibilities into a choice from $n$ possibilities with probabilities $p_1, \dots, p_n$ and then, if the $i$th possibility was chosen, $n_i$ choices of equal probability $p_i.$ We do this because above, we found how to find $H$ when all events are equally likely, and property (3) of our desired function lets us break down our overall choice from $\sum n_i$ possibilities. This gives us

\bal
K \log{\sum n_i} & = H(p_1, \dots, p_n) + \sum p_i A(n_i) \\
                 & = H(p_1, \dots, p_n) + \sum p_i K \log{n_i} \\
                 & = H(p_1, \dots, p_n) + K \sum p_i \log{n_i}. \\
\eal

Then,
\bal
H(p_1, \dots, p_n) & = K \log{\sum n_i} - K \sum p_i \log{n_i}  \\
                   & = K \left [ \log{\sum n_i} - \sum p_i \log{n_i}  \right ] \\
                   & = K \left [ \sum p_i \left ( \log{\sum n_i} \right ) - \sum p_i \left ( \log{n_i} \right )  \right ] \quad \text{ because } \sum p_i = 1  \\
                   & = K \left [ \sum p_i \left ( \log{\sum n_i} -  \log{n_i} \right )  \right ]  \\
                   & = K \left [ \sum p_i \left ( - \log{\frac{n_i}{\sum n_i}} \right )  \right ]  \\
                   & = - K  \sum p_i \log{p_i }.  \\
\eal

If the $p_i$ are incommensurable, we can approximate them as closely as we'd like with \@{rationals}, since \@{rationals-are-dense-in-reals}. By the first property we assumed for $H,$ it is \@{continuous} in the $p_i,$ and so its value at the incommensurable $p_i$ equals its \@{limit} as we approach via the \@{rationals}, and so our expression holds in general. $K$ is left to us to pick, picking it is equivalent to picking a base for the logarithm.
\end{proof}

\detach

\begin{remark}
The \@{entropy} of a \@{random-variable} quantifies the average level of uncertainty or information associated with the variable's possible outcomes. It measures the expected amount of information needed to describe the state of the variable, considering the distribution of probabilities across all potential states.
\end{remark}

The demo below makes these quantities concrete. It uses a discrete alphabet of $N$ symbols whose probabilities follow a \emph{Zipfian} distribution,

\[ p(n) = \frac{n^{-\alpha}}{\sum_{k=1}^{N} k^{-\alpha}}, \qquad n = 1, 2, \dots, N, \]

where the skew $\alpha$ controls how unevenly probability mass is spread across symbols. At $\alpha = 0$ every symbol is equally likely (the uniform distribution), so \@{entropy} is maximal at $\log N.$ As $\alpha$ grows the distribution concentrates on the first few symbols, the rare symbols carry ever more \@{self-information}, and the \@{entropy} falls toward $0.$ The bar chart shows the probability of each symbol alongside its \@{self-information}; the second chart traces the \@{entropy} as $\alpha$ sweeps from uniform to highly skewed.

\includedemo{zipf-entropy}

Examples to cover: [https://claude.ai/chat/0702d3b2-a181-4b5b-a572-e5464d9d24d5]

\subsection{Example: Coin Flip}

For a fair coin flip, $\mathcal{X} = \{h, t\}$ and $p(h) = p(t) = 0.5.$ Then, $I(h) = I(t) = - \log{(0.5)} = 1 \text{ bit},$ and

\[ H(X) = p(h) \cdot I(h) + p(t) \cdot I(t) = 0.5 \cdot 1 + 0.5 \cdot 1 = 1 \text{ bit}. \]

Let's say we have an unfair coin and $p(h) = 0.9, p(t) = 0.1.$ Then

\[ I(h) = - \log{(0.9)} \approx 0.152 \text{ bits}, \quad I(t) = - \log{(0.1)} \approx 3.32 \text{ bits}, \]

and for \@{entropy} of $X$ we get

\[ H(X) = p(h) \cdot I(h) + p(t) \cdot I(t) \approx 0.9 \cdot 0.152 + 0.1 \cdot 3.32 \approx 0.469 \text{ bits}. \]

\subsection{Example: Six-Sided Die}

For a fair six-sided die, $\mathcal{X} = \{1,2,3,4,5,6\}$ and $p(x) = 1/6$ for all $x.$ Then, $I(x) = - \log{(1/6)} = \log(6) \approx 2.58 \text{ bits},$ and

\[ H(X) = 6 \cdot \frac{1}{6} \cdot I(x) = I(x) \approx 2.58 \text{ bits}. \]

Let's say we have an unfair die and $p(1) = 0.5$ and $p(x) = 0.1$ for $x \in \{2,3,4,5,6\}.$ Then

\[ I(1) = - \log{(0.5)} = 1 \text{ bit}, \quad I(x) = - \log{(0.1)} \approx 3.32 \text{ bits for } x \neq 1, \]

and for \@{entropy} of $X$ we get

\[ H(X) = p(1) \cdot I(1) + \sum_{x=2}^{6} p(x) \cdot I(x) \approx 0.5 \cdot 1 + 5 \cdot 0.1 \cdot 3.32 \approx 2.16 \text{ bits}. \]

\subsection{Some Theorems on Entropy}

\begin{theorem}
When a \@{random-variable} is uniformly distributed over an alphabet of $n$ \@{elements}, the \@{self-information} of any given \@{element} equals the \@{entropy} of the \@{random-variable} and is $\log{n}.$
\end{theorem}

\begin{theorem}
$H = 0$ iff all the $p_i$ but one are zero, this one having the value of one.
\end{theorem}

\begin{proof}
Suppose $H(p_1, \dots, p_n) = 0.$ Then, $0 = - \sum_{i=1}^n p_i \log{p_i}.$ Note that because $0 \leq p_i \leq 1,$ we have $\log{p_i} \leq 0,$ and $-p_i \log{p_i} \geq 0$ (with the convention $0 \log 0 = 0,$ since $\lim_{p \to 0^+} p \log p = 0$).  Assume for contradiction that more than one $p_i$ is non-zero. Then, because $\sum p_i = 1,$ each non-zero $p_i$ is in $(0, 1)$ and therefore its $- p_i \log{p_i} > 0,$ and the sum of these terms is therefore non-zero, a contradiction. Now, since all probabilities are required to sum to 1, it can't be the case that all probabilities are zero, which means that exactly one probability must be non-zero and that probability must be $1.$
\end{proof}

\begin{theorem}
A \@{random-variable} has the most \@{entropy} when the \@{elements} of its alphabet are all equally likely to occur.
\end{theorem}

\begin{proof}
We will use a \@{Lagrangian} function of \@{entropy} to show this. Let

\[ \mathcal{L}(\vec{p}, \lambda) = - \sum_i \left [ p_i \ln{p_i} \right ] - \lambda \left [ \left ( \sum_i p_i \right ) - 1 \right ]. \]

Then, for any $p_i$ we want

\bal 
\frac{\partial \mathcal{L}}{\partial p_i} & = -\ln{p_i} - 1 - \lambda = 0 \\
                                          & \implies p_i = e^{-1 - \lambda}.
\eal

But, this means $p_i$ doesn't depend on $i,$ and so is the same for each $i$ and therefore $p_i = \frac{1}{n}.$

It remains to show that this extreme of $H$ is a maximum. This follows from the facts that the domain is a \@{convex}, \@{compact} set and that the entropy function is strictly \@{concave}. TODO: more details on these.
\end{proof}

\begin{definition}[Joint Entropy]
	The \textbf{joint entropy} $H(X,Y)$ of a pair of \@{discrete} \@{random-variables} $(X,Y)$ with a \@{joint-distribution} $p(x,y)$ is defined as

\[ H(X,Y) = -\sum_{x \in \mathcal{X}} -\sum_{y \in \mathcal{y}} p(x,y) \log{p(x,y)}, \]

which can also be expressed as

	\[ H(X,Y) = - \mathbb{E} \left [ \log{p(X,Y)} \right ]. \]

\end{definition}


\begin{theorem}\label{joint-entropy-is-less-than-or-equal-to-entropy-of-parts}
Let $X$ and $Y$ be discrete \@{random-variables} with alphabets $\mathcal{X}$ and $\mathcal{Y},$ and let $p(x,y)$ be the probability of the joint occurrence of $X = x$ and $Y = y.$

	The \@{entropy} of the joint event (the \@{joint-entropy}), $H(X,Y),$ is less than or equal to the sum of the individual entropies, i.e.

\[ H(X,Y) \leq H(X) + H(Y), \]

with equality iff $X$ and $Y$ are independent, that is, iff $p(x,y) = p(x)p(y).$
\end{theorem}

\begin{proof}
The \@{entropy} of the joint event is

\[ H(X,Y) = - \sum_{x,y} p(x,y) \log{p(x,y)}. \]

This is just treating each possible pair of individual outcomes as its own outcome, i.e. we have $|\mathcal{X}| \times |\mathcal{Y}|$ possible outcomes. Writing $p(x) = \sum_{y'} p(x,y')$ and $p(y) = \sum_{x'} p(x',y)$ for the marginals (the argument names the distribution), we have

\[ H(X) = - \sum_{x,y} p(x,y) \log{\sum_{y'}{p(x,y')}}, \quad H(Y) = - \sum_{x,y} p(x,y) \log{\sum_{x'}{p(x',y)}}. \]

So,

\bal
H(X) + H(Y) - H(X,Y) & = - \sum_{x,y} p(x,y) \log{\sum_{y'}{p(x,y')}} - \sum_{x,y} p(x,y) \log{\sum_{x'}{p(x',y)}} + \sum_{x,y} p(x,y) \log{p(x,y)} \\
                    & = \sum_{x,y} p(x, y) \left ( \log{p(x,y)} - \log{\sum_{y'}{p(x,y')}} - \log{\sum_{x'}{p(x',y)}}  \right ) \\
                    & = \sum_{x,y} p(x, y) \left ( \log{p(x,y)} - \log{p(x)} - \log{p(y)}  \right ) \\
                    & = \sum_{x,y} p(x, y) \left ( \log{\frac{p(x,y)}{p(x)p(y)}}  \right ) \\
                    & = - \sum_{x,y} p(x, y) \left ( \log{\frac{p(x)p(y)}{p(x,y)}}  \right ) 
\eal

Now, let's define a new random variable for convenience, $Z = \frac{p(X)p(Y)}{p(X,Y)}.$ We now have that

\bal
H(X) + H(Y) - H(X,Y) & = - \sum_{x,y} p(x, y) \left ( \log{\frac{p(x)p(y)}{p(x,y)}}  \right ) \\
                     & = - \sum_{x,y} p(x, y) \left ( \log{Z}  \right ) \\
                     & = -\mathbb{E}[\log{Z}]. \\
\eal

Now, by \@{jensens-inequality}, because \@{log-is-concave}, we have that $\mathbb{E}[\log{Z}] \leq \log{\mathbb{E}[Z]}.$ Now,

\bal
\mathbb{E}[Z] & = \sum_{x,y} p(x,y) \cdot \frac{p(x)p(y)}{p(x,y)} \\
     & = \sum_{x,y}p(x)p(y) \\
     & = \big ( \sum_{x}p(x) \big ) \big ( \sum_{y}p(y) \big )  \\
     & = 1 \cdot 1 \\
     & = 1. \\
\eal
Therefore $\log{\mathbb{E}[Z]} = \log{1} = 0,$ and we have that $\mathbb{E}[\log{Z}] \leq 0 \implies -\mathbb{E}[\log{Z}] \geq 0,$ and

\bal
H(X) + H(Y) - H(X,Y) & \geq 0 \\
- H(X,Y) & \geq -H(X) -H(Y) \\
H(X,Y) & \leq H(X) + H(Y). \\
\eal

Now, for the equality part. Suppose $H(X) + H(Y) = H(X,Y).$ Then, $\mathbb{E}[\log{Z}] = 0,$ and recalling that $\mathbb{E}[Z] = 1,$ $\mathbb{E}[\log{Z}] = \log{\mathbb{E}[Z]} = 0.$ In \@{jensens-inequality}, equality holds only if our function is \@{affine} or if $Z$ is constant; since $\log$ is \@{affine} on no interval, $Z$ must be constant, and since $\mathbb{E}[Z] = 1,$ $Z = \frac{p(x)p(y)}{p(x,y)} = 1 \implies p(x)p(y) = p(x,y).$ 

Now, suppose $p(x)p(y) = p(x,y).$ Then $Z = 1$ and

\[ H(X) + H(Y) - H(X,Y) = - \sum_{x,y} p(x, y) \left ( \log{1}  \right ) = 0, \]

so $H(X) + H(Y) = H(X,Y).$
\end{proof}

\detach

\begin{note}
When we decrease the probability of an \@{element} occurring, its \@{self-information} increases only logarithmically, while its weight in the entropy sum decreases linearly.
\end{note}

\begin{definition}[Cross-Entropy]
Given two discrete probability distributions, $p$ and $q,$ that share a support $\mathcal{X},$ \textbf{cross-entropy} measures the average number of bits needed to represent an event drawn from $\mathcal{X}$ when the coding scheme is optimized for an estimated distribution $q$ rather than the true distribution $p.$ Its formula is very similar to that of \@{entropy}, except \@{surprisal} is calculated using $q$ while \@{expectation} uses $p$:

\[ H(p,q) = - \sum_{x \in \mathcal{X}} p(x) \log{q(x)}. \]

We can write it in \@{expectation} form as

\[ H(p,q) = \mathbb{E}_p \left [- \log{q}  \right ], \]

where $\mathbb{E}_p$ is expectation under $p.$
\end{definition}

Note that when $p = q,$ \@{cross-entropy} equals \@{entropy}. The symbol $H$ is overloaded: $H(p,q)$ with distributions as arguments means \@{cross-entropy}, while $H(X,Y)$ with \@{random-variables} as arguments means joint entropy.

\begin{definition}[KL Divergence]\synonyms{relative-entropy, "Kullback-Leibler divergence"}
For discrete probability distributions, $p$ and $q,$ that share a support $\mathcal{X},$ the \textbf{relative entropy} from $q$ to $p$ is defined to be

\[ D_{KL}(p || q) = \sum_{x \in \mathcal{X}} p(x) \log{\frac{p(x)}{q(x)}}, \]

which is just \@{cross-entropy} minus \@{entropy}:

\[ D(p || q) = H(p,q) - H(p). \]

Note that we don't always use the $KL$ subscript where it's obvious from the context.

We can also write it as

\[ D(p || q) = \mathbb{E}_p \left [ - \log{q} \right ] - \mathbb{E}_p \left [ - \log{p} \right ]  =  \mathbb{E}_p \left [ \log{\frac{p}{q}}  \right ], \]

which makes it clear that it's the expected inefficiency in an encoding optimized for $q$ rather than $p.$
\end{definition}

\begin{remark}
\@{kl-divergence} tells us how different an approximating distribution $q$ is from a true probability distribution $p.$ It's not a true distance metric - it's asymmetric, but it does behave like distance in that it gets smaller as $q$ gets more like $p$ and is zero when $q = p.$ It is a sort of directed distance.

Another way to think about it is as a measure of how inefficient a coding scheme optimized for $q$ is when the actual distribution is $p.$ If $q$ says some event is rare, and so we use a longer symbol for it (more bits), but the event actually occurs frequently, we'll waste bits representing the event when we could have used a shorter symbol for it.
\end{remark}

\begin{theorem}[Gibbs' Inequality]\label{gibbs-inequality}
For discrete probability distributions $p$ and $q$ that share a support $\mathcal{X},$

\[ D(p || q) \geq 0, \]

with equality iff $p = q.$ Equivalently, \@{cross-entropy} is never less than \@{entropy}: $H(p,q) \geq H(p),$ with equality iff $q = p.$
\end{theorem}

\begin{proof}
Because \@{log-is-concave}, by \@{jensens-inequality} we have that

\bal
\mathbb{E}_p \left [ \log{\frac{q_i}{p_i}}  \right ] & \leq \log{ \left ( \mathbb{E}_p \left [ \frac{q_i}{p_i} \right ] \right ) }  \\
                                                     & = \log{ \left ( \sum_i \left [ p_i \frac{q_i}{p_i} \right ] \right ) }  \\
                                                     & = \log{ \left ( \sum_i q_i \right ) }  \\
                                                     & = \log{ 1 }  \\
                                                     & = 0.  \\
\eal

Then, since $\mathbb{E}_p \left [ \log{\frac{q_i}{p_i}}  \right ] = -D(p || q),$ we have that $D(p || q) \geq 0.$

Now, note that if $p = q,$ $\frac{q_i}{p_i} = 1$ and $D(p || q) = 0.$ Conversely, suppose $D(p || q) = 0.$ Since $\log$ is strictly \@{concave}, it is not \@{affine}, and \@{jensens-inequality} condition for equality gives us that $\frac{q_i}{p_i} = c.$ Now, $\sum_i q_i = \sum_i c p_i = 1,$ and since $p$ is a probability distribution, $c = 1,$ and therefore $p_i = q_i$ and $p = q.$
\end{proof}

\begin{theorem}[Doubly stochastic maps increase entropy]
Take a \@{doubly-stochastic} \@{matrix} $A$ and \@{probability-vector} $p,$ let $q = A p.$ Then $H(q) \geq H(p),$ with equality iff $q$ is a \@{rearrangement} of $p$.
\end{theorem}

\begin{proof}
Note that $q_i = \sum_{j} A_{ij} p_j,$ by the definition of vector matrix multiplication. Now we'll define a couple of joint distributions:

\[ r_{ij} = A_{ij} p_j, \quad s_{ij} = A_{ij} q_i. \]

Now we find the \@{relative-entropy} from $s$ to $r$:

\bal
D(r || s) & = \sum_{i,j} r_{ij} \log{\frac{r_{ij}}{s_{ij}}} \\
          & = \sum_{i,j} A_{ij} p_j \log{\frac{A_{ij} p_j}{A_{ij} q_i }} \\
          & = \sum_{i,j} A_{ij} p_j \log{\frac{p_j}{q_i}} \\
          & = \sum_{i,j} A_{ij} p_j \left ( \log{p_j} - \log{q_i} \right ) \\
          & = \sum_{i,j} A_{ij} p_j \log{p_j} - \sum_{i,j} A_{ij} p_j \log{q_i} \\
\eal

Now, addressing the first term on the RHS of the last line:

\bal
\sum_{i,j} A_{ij} p_j \log{p_j} & = \sum_j \left [ p_j \log{p_j} \sum_{i} A_{ij} \right ] \\
                                & = \sum_j \left [ p_j \log{p_j} \cdot 1 \right ] \\
                                & = \sum_j  p_j \log{p_j} \\
                                & = -H(p). \\
\eal

And the second term on the RHS of the last line:

\bal
\sum_{i,j} A_{ij} p_j \log{q_i} & = \sum_{i} \left [ \log{q_i} \sum_{j} A_{ij} p_j \right ] \\
                                & = \sum_{i} \left [ \log{q_i} q_i \right ] \\
                                & = -H(q).
\eal

So, we end up with $D(r || s) = H(q) - H(p).$ Since $D(r || s) \geq 0$ by \@{gibbs-inequality}, we have $H(q) - H(p) \geq 0,$ which implies that $H(q) \geq H(p),$ which is what we wanted to show.

Now, the equivalence case. If $q$ is just a rearrangement of $p,$ then obviously it has the same \@{entropy} as we can just re-index to recover $p.$ Now assume $H(p) = H(q).$ Then, $H(q) - H(p) = 0 \implies D(r || s) = 0 \implies r = s \implies A_{ij} p_j = A_{ij} q_i \implies p_j = q_i$ whenever $A_{ij} > 0.$ Now, for each value $c$ appearing in $p$ or $q,$ let

\[ J_c = \{j : p_j = c\}, \quad I_c = \{i : q_i = c\}, \]

that is $J_c$ is the set of indices of $p$ where $p_j = c,$ and similarly with $I_c$ and $q.$ Now, consider row $i \in I_c.$ For any $A_{ij} > 0,$ we have $p_j = q_i,$ so $j \in J_c,$ i.e. every nonzero entry of row $i$ lies in a column of $J_c,$ so the columns $j \in J_c$ have the entire mass of row $i.$ Similarly, for column $j \in J_c,$ for nonzero $A_{ij},$ $q_i = p_j = c,$ so $i \in I_c,$ and the rows $I_c$ contain all of column $j$'s mass. Taking $\vec{1}_S$ as the \@{indicator-vector} of a set $S,$ we have that

\[ A \vec{1}_{J_c} = \vec{1}_{I_c}, \]

that is, rows in $I_c$ sum to $1$ over the columns $J_c,$ while rows outside $I_c$ have no mass in the columns $J_c$ at all. Now we count the mass:

\bal
|I_c| & = \vec{1}^\top \vec{1}_{I_c}  \qquad \text{count entries via dot product} \\
      & = \vec{1}^\top (A \vec{1}_{J_c}) \qquad \text {identity from above} \\
      & = (\vec{1}^\top A) \vec{1}_{J_c} \qquad \text{associativity}  \\
      & = \vec{1}^\top \vec{1}_{J_c} \qquad \text{column-stochasticity}  \\
      & = |J_c| \qquad \text{count entries via dot product}.
\eal

Therefore, $|I_c| = |J_c|$ for all $c,$ so each $c$ appears equally often in $p$ and in $q,$ and $q$ is therefore a rearrangement of $p.$
\end{proof}

\begin{intuition}
Averaging (well, replacing each probability with a \@{convex-combination} of all the probabilities in a way that retains their summing to 1) the probabilities in $p$ can only bring them closer to each other (if it doesn't just rearrange them), which increases \@{entropy}.
\end{intuition}

\begin{definition}[Conditional Entropy]
Let $X$ and $Y$ be \@{random-variables}, not necessarily \@{independent}. The probability that $Y$ takes the value $y$ when $X$ takes the value $x$ (for any specific $x \in \mathcal{X}, y \in \mathcal{Y}$) is called the \@{conditional-probability} $p(y|x)$ and is given by:

\[ p(y|x) = \frac{p(x, y)}{p(x)}, \quad p(x) = \sum_{y} p(x,y). \]

(Note that $p(x,y)$ is the \@{joint-probability} of $X = x, Y = y.$)

We define the \textbf{conditional entropy} of $Y$ given $X,$ $H(Y|X),$ as the average of the entropy of $Y$ for each value $x$ of $X$ weighted according to the probability of getting that particular $x.$ That is,

\[ H(Y|X) = -\sum_{x \in \mathcal{X}, y \in \mathcal{Y}} p(x,y) \log{p(y|x)}. \]
\end{definition}

\begin{note}
We can obtain the above formula for \@{conditional-entropy} as follows. First, if $X = x,$ then our \@{entropy} for $Y$ is

\[ H(Y | X = x) = -\sum_{y} p(y|x) \log{p(y|x)}. \]

Now, summing across $X$ gives:

\bal
H(Y|X) & = \sum_{x} p(x) H(Y | X = x) \\
       & = \sum_{x} p(x) \left [  -\sum_{y} p(y|x) \log{p(y|x)}  \right ]  \\
       & = - \sum_{x} \sum_{y}\left [ p(x)  p(y|x) \log{p(y|x)}  \right ]  \\
       & = - \sum_{x, y} p(x,y) \log{p(y|x)}. \\
\eal
\end{note}

\begin{theorem}[Chain rule for joint entropy]\label{joint-entropy-is-base-plus-conditional}
	The \@{entropy} (\@{joint-entropy}) of the joint event $X,Y$ is the \@{entropy} of $X$ plus the \@{entropy} of $Y$ when $X$ is known.
\end{theorem}

\begin{proof}
Note that

\bal
H(Y | X) & = - \sum_{x, y} p(x,y) \log{\frac{p(x,y)}{p(x)}} \\
         & = - \sum_{x, y} p(x,y) \log{p(x,y)} + \sum_{x, y} p(x,y) \log{p(x)}.\\
\eal

The first term on the RHS is just $- \sum_{x, y} p(x,y) \log{p(x,y)} = H(X,Y).$ The second is

\bal
\sum_{x, y} p(x,y) \log{p(x)} & = \sum_{x} \left [ \sum_{y} p(x,y) \right ] \log{p(x)}  \\
                              & = \sum_{x} p(x) \log{p(x)}  \\
                              & = -H(X).
\eal

Altogether, we have that $H(Y | X) = H(X,Y) - H(X),$ or

\[ H(X,Y) = H(X) + H(Y | X). \]
\end{proof}

\begin{theorem}
The \@{uncertainty} of $Y$ is never increased by knowledge of $X.$ It will be decreased unless $X$ and $Y$ are independent events, in which case it is not changed.
\end{theorem}

\begin{proof}
From \@{joint-entropy-is-less-than-or-equal-to-entropy-of-parts} and \@{joint-entropy-is-base-plus-conditional}, we have

\bal
H(X) + H(Y) & \geq H(X,Y) \\
            & = H(X) + H(Y|X), \\ 
\eal

hence $H(Y) \geq H(Y|X).$
\end{proof}

\end{document}
